\emph{Pat} \cite{fowlerSpecialDeliveryProgramming2023} is the first programming language to support mailbox types.
In contrast to languages like Erlang or Elixir, where mailboxes are implicit,
Pat supports first-class explicit mailboxes. 
We can express the Future example in Pat as follows:

\begin{minipage}{0.59\textwidth}
\begin{pat}[firstnumber=1]
def $\textsf{emptyFuture}$($\mathit{self} : \mathsf{EmptyFuture}$) : $\unittype$ {
  guard $\mathit{self}$ : $\msg{Put} \mbcomp \mbstar{\msg{Get}}$ {
      receive $\msg{Put}$[$x$] from $\mathit{self}$ |->
      $\mathsf{fullFuture}(\mathit{self}, x)$
  }
}
\end{pat}
\begin{pat}[firstnumber=7]
def $\textsf{fullFuture}$($\mathit{self} : \mathsf{FullFuture}, \mathit{value} : \mathsf{Int}$) : $\unittype$ {
  guard $\mathit{self}$ : $\mbstar{\msg{Get}}$ {
    free |-> ()
    receive $\msg{Get}$[$\mathit{user}$] from $\mathit{self}$ |->
        $\mathit{user}\mathsf{!}\msg{Reply}[value]$;
      $\mathsf{fullFuture}(\mathit{self}, \mathit{value})$
  }
}
\end{pat}
\end{minipage}
\begin{minipage}{0.5\textwidth}
\begin{pat}[firstnumber=15]
def $\textsf{client}$() : $\unittype$ {
  $\keyword{let}$ $\var{future}$ = $\TNew$ in
  spawn $\textsf{emptyFuture}$($\var{future}$);
  let $\mathit{self}$ = $\TNew$ in
  $\TSend{\var{future}}{Put}{5}$;
  $\TSend{\var{future}}{Get}{\mathit{self}}$;
  guard $\mathit{self}$ : $\msg{Reply}$ {
      receive $\msg{Reply}$[$\mathit{result}$]
    from $\mathit{self}$ |->
      free $\mathit{self}$;
      $\mathsf{print}(\mathsf{intToString}(\mathit{result}))$
  }
}
\end{pat}
\end{minipage}
\vspace{.2cm}

To support dynamic mailbox creation, new mailboxes are created using the
$\keyword{new}$ keyword and bound to names using $\keyword{let}$-bindings.
Pat programs can then send to mailboxes (e.g., $\TSend{\var{future}}{Put}{5}$
sends a $\msg{Put}$ message with payload $5$ to mailbox \var{future}) and
receive from mailboxes using the $\keyword{guard}$ keyword.
A guard is always accompanied by a mailbox pattern, representing the mailbox's
contents, and contains \emph{guard clauses} including $\keyword{receive}$ which
receives a message from the mailbox and binds the payload and new mailbox name;
and $\keyword{free}$ which is invoked when a mailbox will not receive any more
messages and can be deallocated.

%Since mailbox names are first-class, it is possible to send some mailbox name $x$ to a function that contained a mailbox name $x$ in its scope, but is already free in that context.
%This would result in a \emph{use-after-free} error.
% The authors identify three properties needed in every process to avoid these problems \cite{fowlerSpecialDeliveryProgramming2023}:
% \begin{enumerate}
%   \item No two distinct variables should represent the same underlying mailbox name.
%   \item Once let-bound to a different name, a mailbox variable is considered out-of-scope.
%   \item A mailbox name cannot be used after it has been used in a guard.
% \end{enumerate}


A mechanization of the operational semantics of Pat is outside of the scope of
this paper, but we give an outline of the semantics here for context.
Pat's semantics makes use of a \emph{frame stack}
representation that simplifies specifying inductive invariants for the type
preservation proof.
There are two reduction relations: a deterministic
$\beta$-reduction relation on functional terms (like $\keyword{let}$-bindings
and function calls) that is entirely standard, and a nondeterministic reduction
relation on a language of \emph{configurations} that closely resemble processes
in the $\pi$-calculus. The communication and concurrency constructs reduce as
follows:

\begin{itemize}
    \item Evaluating the $\TNew$ construct creates a fresh runtime mailbox name and returns
        it to the calling thread.
    \item Evaluating a message send $\TSend{a}{m}{\vala}$ returns the unit value
        to the calling thread and creates an \emph{in-flight} message
        $\sendconfig{a}{m}{\vala}$.
    \item Evaluating $\TSpawn{\tma}$ returns the unit value to the calling
        thread and evaluates $\tma$ in a new process.
    \item The behaviour of a $\TGuard{a}{E}{\seq{\gd}}$ expression depends on
        the messages that have been sent to mailbox $a$ along with the guards
        contained in $\seq{\gd}$:
        \begin{itemize}
            \item If there are no messages sent to $a$; no references to $a$
                in other threads; and $\seq{\gd}$ contains a $\GFree{\tma}$
                guard, then mailbox name $a$ will be deleted and the thread will
                run computation $\tma$.
            \item If there is a message $\sendconfig{a}{m}{\vala}$ and
                $\seq{\gd}$ contains a clause $\GReceive{m}{\tma}$, then
                the thread will run computation $\tma$ with bindings for the
                received value and updated mailbox name.
    \end{itemize}
\end{itemize}


A major problem in Pat is \emph{mailbox name aliasing}, where multiple static
names could be given to the same underlying mailbox name. This could be used to
break type soundness.
The Mailbox Calculus addresses this issue using a dependency graph, which can
additionally rule out cyclic dependencies, but this approach does not scale to
programming languages (see~\cite{fowlerSpecialDeliveryProgramming2023} for a detailed discussion).
Instead, Pat uses \emph{quasi-linear typing}
\cite{kobayashiQuasilinearTypes1999,ennalsLinearTypesPacket2004a}, where types
are equipped with \emph{usage annotations}.  If a type is defined as
\emph{returnable}, it is allowed to appear in the return type of an expression
and be returned from an evaluation frame.
Otherwise, if a type is defined as \emph{second-class},
it can be used several times, but is never allowed to escape its scope.
%
A returnable occurrence of a mailbox name must be its last lexical usage.
%
Quasilinearity annotations, in combination with syntactic restrictions,
eliminates aliasing issues and self-deadlocks.

The authors of Pat present both a \emph{declarative} and an
\emph{algorithmic} version of their type system.  The algorithmic version (which is needed because the
declarative system makes use of environment subtyping and non-deterministic
context splits) makes use of
\emph{bidirectional typing} and \emph{constraint solving}, and is at the core of
Pat's implementation.  This paper focuses on the declarative type
system.
